%% research-memo.tex — Tufte-style research memo template
\documentclass[letterpaper, justified, notitlepage]{tufte-handout}

%% ── Packages ────────────────────────────────────────────────────────────────
\usepackage[T1]{fontenc}
\usepackage{palatino}
\usepackage{microtype}
\usepackage{booktabs}
\usepackage[table]{xcolor}
\usepackage{graphicx}
\usepackage{amsmath, amssymb}
\usepackage{amsthm}
\usepackage{hyperref}
\usepackage{enumitem}
\usepackage[most]{tcolorbox}

%% ── Theorem environments ────────────────────────────────────────────────────
\theoremstyle{plain}
\newtheorem{definition}{Definition}[section]

%% ── Typography ──────────────────────────────────────────────────────────────
\setlist{noitemsep, topsep=2pt}
\setcounter{secnumdepth}{2}   % number sections and subsections

%% ── Memo fields (fill these in) ─────────────────────────────────────────────
\newcommand{\memoto}{%
  Jennifer, Thomas, Claire
}
\newcommand{\memofrom}{%
  Kyle
}
\newcommand{\memodate}{\today}
\newcommand{\memosubject}{%
  Notes on the empirical methods used in related work
}
\newcommand{\memosummary}{%
}

%% ── Callout box ─────────────────────────────────────────────────────────────
\definecolor{calloutblue}{RGB}{219,234,254}
\newcommand{\callout}[1]{%
  \begin{tcolorbox}[
    colback=calloutblue,
    colframe=black,
    arc=6pt,
    boxrule=0.4pt,
    left=6pt, right=6pt, top=4pt, bottom=4pt
  ]
    #1
  \end{tcolorbox}%
}

%% ── Printer's mark box ──────────────────────────────────────────────────────
\definecolor{markgray}{gray}{0.88}

%% ────────────────────────────────────────────────────────────────────────────
\begin{document}

%% Full-width memo header
\begin{fullwidth}
  \noindent
  \begin{minipage}[t]{0.62\linewidth}
    \vspace{0pt}
    \begin{tabular}{@{}lp{0.72\linewidth}}
      \textbf{To:}      & \memoto      \\[1pt]
      \textbf{From:}    & \memofrom    \\[1pt]
      \textbf{Date:}    & \memodate    \\[1pt]
      \textbf{Subject:} & \memosubject \\
    \end{tabular}
  \end{minipage}%
  \hfill
  \begin{minipage}[t]{0.22\linewidth}
    \vspace{0pt}
    \raggedleft
    \colorbox{markgray}{%
      \parbox[c][2.2cm][c]{0.95\linewidth}{%
        \centering\color{gray}\footnotesize [printer's mark]
      }%
    }
  \end{minipage}

  \vspace{0.6em}
  \noindent\rule{\linewidth}{0.5pt}
  \vspace{0.6em}

  \noindent\textbf{Summary:}\quad\memosummary

  \vspace{0.8em}
  \noindent\rule{\linewidth}{0.4pt}
  \vspace{1em}
\end{fullwidth}

%% ────────────────────────────────────────────────────────────────────────────
%% Body — single column with wide right margin for notes and figures
%% Use \sidenote{} for numbered margin notes (replaces \footnote)
%% Use \marginnote{} for unnumbered margin notes
%% Use \begin{marginfigure}...\end{marginfigure} for margin figures
%% ────────────────────────────────────────────────────────────────────────────

\section{Bailey and Rottinghaus (2014)}
Extract the enabling clause from dataset composed of proclamations and executive orders. Then,
they label each document according to the source of the authority in the enabling clause. They use 
three labels: President-based, Congress-based, or joint. President-based authorities include an
invocation of President's powers as president generally, or their powers under the Constitution.
Congress-based are signaled by citation of a joint resolution of Congress, specific public law,
specific statute, specific U.S. Code, or declaration of "laws of the United States." Documents
labelled "joint" include both a President- and Congress-based authority, but are treated as
Congress-based for the analysis.\sidenote{See footnote 5 of the paper.}

Documents are hand-labelled using a single annotator, then evaluated for consistency by a second
annotator.\sidenote{See footnote 4 of the paper.}

\section{Kaufman and Rogowski (2024); Prez policymaking}
\textbf{Questions}: (1) Is a document significant (binary) and (2) What issue area is the document?

Both are document-level, features are bag of words (either unigram vs. bigram).


\section{Kaufman and Rogowski (2024); Strategic Substitution}
\textbf{Question}: Do presidents substitute less visible types of directives for EOs or
proclamations under divided govt? Goal to avoid Congressional scrutiny.

Dataset composed of 33,921 directives from 1946-2020. Includes memoranda, public land orders,
executive agreements, agency directives, EOs, and proclamations. 

No classification, just use directives labelled significant from the \textit{Presidential
policymaking} paper. Then use linear probability model to test whether possible to predict type of
directive from whether government divided, controlling for stuff. 

\section{Ruhl et al (2018)}
Uses LDA and BoW to topic model executive orders, presidential memoranda, proclamations, and
determinations according to the document's content. Examples of labels:
\begin{itemize}
  \item Military positions, succession, titles
  \item Internal emergencies and sanctions.
\end{itemize}




\end{document}
